% !TEX root = ../main.tex
\clearpage
\section{Adquisición comercial y desarrollo propio: consideraciones técnicas}
\label{app:construir}
El análisis de una alternativa de desarrollo propio comprende componentes de procesamiento, recursos de evaluación y requisitos de integración. La documentación disponible incluye los modelos multimodales de pesos abiertos Gemma 4 y Qwen3-Omni, así como los recursos de evaluación PriMock57 y TEME-v1~\cite{J4955P47,P9P5U4KE,G2L7QBIV,CMRQCQTY}. Estos elementos pueden formar parte de una solución, pero su disponibilidad no equivale a disponer de un escriba clínico validado, certificado e integrado con HCIS.

Y-KNOT ilustra la normalización de documentos, la definición de eventos de activación y la devolución de borradores para aprobación dentro de la historia clínica electrónica. Su entrada se compone de registros existentes y el código no es público; estas características limitan su consideración como alternativa abierta a la captura ambiental~\cite{QQCRZ7VX}.

La comparación técnica debe incluir licencias y versiones, procesamiento local o remoto, capacidad para audio en español, trazabilidad de errores, conectores, autenticación, mantenimiento, actualización y coste total. Tanto una arquitectura de reconocimiento automático del habla y generación de texto (ASR+LLM) como un modelo multimodal requieren evaluar la nota final y el flujo completo. La selección de información y la revisión clínica siguen siendo funciones necesarias, con independencia de su distribución entre componentes.

\pendiente{Este análisis preliminar requiere una selección sistemática de aplicaciones abiertas y una estimación de infraestructura, operación y recursos de desarrollo.}

\printbibliography[heading=subbibliography,title={Referencias técnicas},keyword=tecnica]
